% P2 BSPC V7 evidence-locked comparison; 2026-10-15 author-review ONLY.
\subsection{Informed common-average-reference comparator and source-level null ranks}
To exclude a trivial gain from comparing artifact-sensitive raw waveforms
against an artifact-filtering method, we added a strong fully informed
comparison that subtracts the common average reference (CAR) across
the eight source and target channels prior to applying the unchanged
shared-clock assignment algorithm. The same six original CHB--MIT
EDF digital source prefixes each produced eight controlled reuse
positives and six constructed negatives; these are 84 \emph{correlated}
case variants from six source records, not 84 independent clinical
observations.

On this \emph{controlled} set, raw shared-clock matching detected
48/48 positives with 6/36 false matches; CAR preprocessing followed
by the original shared-clock matcher detected 39/48 with 0/36 false
matches, and the matched-channel residual verification found 48/48
with 0/36 false matches. All six raw false matches occurred in
the deliberately injected strong common-artifact condition.
This result demonstrates a preservation-of-detections advantage
relative to \emph{this specified CAR baseline and its frozen original
thresholds}, but it does not establish advantage over source-calibrated
CAR thresholds, robust rereferencing, component decomposition,
multichannel device synchronization or time-series fingerprints.
The author-review evidence is the inspectable success receipt
\url{https://github.com/lindicaphxag-tech/kaggle/actions/runs/38675387956};
the threshold-tuned CAR check is separately registered
at \url{https://github.com/lindicaphxag-tech/kaggle/actions/workflows/p2-v7-source-tuned-car-benchmark.yml}.
No tuned-CAR result is assumed here.

We also evaluated a continuous score obtained by selecting the
common-lag or one-channel correspondence using a fit-only time
block, followed by verification-block correlation. Rather than
treating individual lags, electrodes or overlapping windows as
independent null trials, an empirical rank test was applied
to \emph{the maximum over the complete permitted source search}
for 19 distinct published CHB--MIT case identifiers. A
controlled positive gave score 0.999968 (common-clock endpoint)
versus maximum null 0.184460; the one-channel endpoint scored
1.000000 versus maximum null 0.424179. Both had the minimum
achievable one-query rank $p=(1+0)/(19+1)=0.05$; the
successful source-level check is
\url{https://github.com/lindicaphxag-tech/kaggle/actions/runs/38675387852}.
The two endpoints are not independent prespecified discoveries;
familywise claims require multiplicity adjustment and suitable
exchangeability of genuinely independent source recordings.
The database is a single corpus and the positives remain engineered
variants, so no population-wide clinical false-match guarantee is
inferred.

\paragraph{Reproducibility and honest scope.}
Before external release, the full numerical witness must be
independently recomputed, the original EDF prefix digests retained
without redistributing patient waveforms, all source-group
identities checked against the distributor manifests, and
threshold selection performed without inspection of final
evaluation recordings. Additional sources and naturally observed
digital reuse, as well as a systematic cost/abstention comparison,
remain required before promoting this component to a generally
superior biomedical signal provenance algorithm.
